\documentclass[11pt]{article}
\usepackage{ochamath}
\subjectcolor{2F5DA8}
\setsheet{線形代数}{固有値・固有ベクトル　演習 1}{https://university-math-crj.pages.dev/linear-algebra/eigenvalues/}
\namelinetrue
\begin{document}
\makesheethead{問題}

\problem[固有値の計算]
次の行列の固有値と固有ベクトルを求めよ。
\[ A = \begin{pmatrix} 2 & 1 \\ 1 & 2 \end{pmatrix} \]
\workspace{50mm}

\problem[対角化]
$A = \begin{pmatrix} 4 & -2 \\ 1 & 1 \end{pmatrix}$ について、$P^{-1}AP$ が対角行列になる正則行列 $P$ を1つ求め、$P^{-1}AP$ を計算せよ。
\workspace{60mm}

\newpage
\problem[連立微分方程式と安定性]
$\dot{\boldsymbol x} = A\boldsymbol x,\quad A = \begin{pmatrix} -1 & 2 \\ 2 & -1 \end{pmatrix},\quad \boldsymbol x(0) = \begin{pmatrix} 2 \\ 0 \end{pmatrix}$ について、
\begin{enumerate}[label=(\arabic*)]
\item 解 $\boldsymbol x(t)$ を求めよ。
\item このシステムは安定か。理由とともに答えよ。
\end{enumerate}
\workspace{60mm}

\problem[2階の微分方程式を行列で書く]
$\ddot y + 3\dot y + 2y = 0$ について、$x_1 = y,\ x_2 = \dot y$ とおく。
\begin{enumerate}[label=(\arabic*)]
\item $\dot{\boldsymbol x} = A\boldsymbol x$ の形に書き、$A$ を求めよ。
\item $A$ の特性方程式が、元の微分方程式の特性方程式 $\lambda^2 + 3\lambda + 2 = 0$ と一致することを示し、固有値を求めよ。
\end{enumerate}
\workspace{40mm}
\end{document}
